\documentclass[letter,11pt]{article}
\usepackage[breaklinks]{hyperref}
\hypersetup{
    bookmarks=true,         % show bookmarks bar?
    unicode=false,          % non-Latin characters in Acrobat’s bookmarks
    pdftoolbar=true,        % show Acrobat’s toolbar?
    pdfmenubar=true,        % show Acrobat’s menu?
    pdffitwindow=false,     % window fit to page when opened
    pdfstartview={XYZ null null 1.00},    % disable zoom
    pdftitle={Homework 1},    % title
    pdfauthor={Richard Zak},     % author
    pdfsubject={UMBC CMSC104 Problem Solving and Computer Programming},   % subject of the document
    pdfkeywords={Computer Science, Programming, Problem Solving, CSEE}, % list of keywords
    pdfnewwindow=true,      % links in new PDF window
    colorlinks=false,       % false: boxed links; true: colored links
    linkcolor=red,          % color of internal links (change box color with linkbordercolor)
    citecolor=green,        % color of links to bibliography
    filecolor=magenta,      % color of file links
    urlcolor=cyan           % color of external links
}
\usepackage{graphicx}
\usepackage{fancyhdr}
\usepackage{multicol}
\pagestyle{fancy}
\usepackage[letterpaper, margin=1in]{geometry}
\geometry{letterpaper}
\usepackage{parskip} % Disable initial indent
\usepackage{color,soul} % Highligher
\usepackage[normalem]{ulem} % Strikethrough with \sout{}

\usepackage[utf8]{inputenc}
\fancyhf{}
\renewcommand{\headrulewidth}{0pt} % Remove default underline from header package
\rhead{CMSC 104 Section 01: Homework 1}
%\rhead{}
\lhead{\begin{picture}(0,0) \put(0,-10){\includegraphics[width=1.1cm]{Images/UMBC-vertical}} \end{picture}}
\cfoot{\thepage}
\rfoot{Spring 2024}
\lfoot{CMSC 104 Section 01}
\AtEndDocument{\vfill \footnotesize{Last modified: 29 August 2022}}
\AtEndDocument{\rfoot{Spring 2024}}
\renewcommand\thesubsection{\arabic{subsection}} % Show only subsection numbers, not section.subsection

\begin{document}

\huge
\textbf{Homework 1: Homework Can Be Fun!}.
\normalsize
\\ ~~ \\
\textbf{Assigned: Thursday 08 September} \\
\textbf{Due Date: Wednesday 14 September}

\section*{Objectives}
\paragraph{}To continue familiarizing with the \texttt{nano} text editor, understanding how to use basic Unix commands (\texttt{cp}, \texttt{mv}, and \texttt{cat}), and reinforcing how to submit assignments. NOTE: Unless otherwise noted, be sure to hit Enter(PC) or return(Mac) after every command.

\section*{Part 1: Original Joke or Motto}
\begin{enumerate}
    \item Login to GL and make sure you are in your home directory (\texttt{pwd}).
    \item Type \texttt{nano joke.txt} to start a new file called ``joke.txt''.
    \item Type your best (clean) joke or personal motto. Edit any mistakes using the \texttt{Backspace} or \texttt{Delete} key to backspace like you would in a normal word processor like Microsoft Word. When you get to the end of a line, hit the enter (PC) or return (Mac) key at a reasonable spot instead of letting the text wrap around to the next line. Remember that \texttt{nano} is simply a text editor and does not format things nicely for us.
    \item Save your work as you go by pressing \texttt{CTRL-O} for ``write out.''
    \item To exit nano, press \texttt{CTRL-X}. If you have made changes since you last saved, it will ask you if you want to save the file, so press \texttt{y} for yes, then pres enter (PC) or return (Mac) when prompted, to save the filename. It is highly advisable NOT to change the filename.
    \item Use the command \texttt{cat joke.txt} to check that your joke or motto was correctly saved.
    \item Have someone check if the joke is funny. If not, \texttt{rm joke.txt} and try another joke or motto.
    \item Submit joke.txt.
    \begin{verbatim}
    [rzak1@linux1 ~]$ submit cmsc104_rzak1 hw1 joke.txt
    \end{verbatim}
    \item Check your submission to make sure the file was properly submitted. If there is no output, \textit{your submission didn’t go through} so repeat step 8, making sure you are in the right directory with the file.
    \begin{verbatim}
    [rzak1@linux1 ~]$ submitls cmsc104_rzak1 hw1
    \end{verbatim}
\end{enumerate}

\section*{Part 2: Welcome to Dad Jokes}
\begin{enumerate}
    \item Make a copy of your, ``joke.txt'', file and call it ``joke2.txt''.
    \begin{verbatim}
    [rzak1@linux1 ~]$ cp joke.txt joke2.txt
    \end{verbatim}
    \item Use \texttt{nano joke2.txt} to edit your ``joke2.txt'' file and add the following joke to the end of the file (or come up with your own!):
    \begin{verbatim}
    Why do ducks have webbed feet?
    To stamp out fires.

    Why do elephants have flat feet?
    To stamp out flaming ducks!
    \end{verbatim}
    \item Save changes to the file by typing \texttt{CTRL-X} to Exit, \texttt{y} to agree to save, and hitting enter (PC) or return (Mac)to save using the provided file name.
    \item Submit your new ``joke2.txt'' file.
    \begin{verbatim}
    [rzak1@linux1 ~]$ submit cmsc104_rzak1 hw1 joke2.txt
    \end{verbatim}
    \item Type logout to log out.
\end{enumerate}

\section*{Grading Rubric}
\begin{itemize}
    \item joke.txt is complete: 45 points
    \item joke2.txt is complete: 45 points
    \item typescript is complete and not garbled: 10 points
\end{itemize}

\section*{What to Submit}
\paragraph{}You should have already submitted the 2 joke files (joke.txt, joke2.txt) by following the instructions above.

\paragraph{}Log back into your GL account. Use the \texttt{script} command to record yourself viewing your joke files using the cat command. Do not record yourself editing your files! Exit from script and submit your typescript file.
\begin{verbatim}
[rzak1@linux1 ~]$ script
Script started, file is typescript
[rzak1@linux1 ~]$ cat joke.txt
NOTE: Your joke/motto should display here.
[rzak1@linux1 ~]$ cat joke2.txt
NOTE: Your joke/motto should display here, followed by the ducks/elephants joke.
[rzak1@linux1 ~]$ exit
exit
Script done, file is typescript
[rzak1@linux1 ~]$ submit cmsc104_rzak1 hw1 typescript
\end{verbatim}

\newpage
\subsection*{Verify}
\paragraph{}Make sure you submitted the assignment correctly.
\begin{verbatim}
[rzak1@linux1 ~]$ submitls cmsc104_rzak1 hw1
\end{verbatim}

\end{document}